% !TeX program = pdflatex
% =====================================================================
%  Rapport de Stage - AI Network Copilot
%  Compatible Overleaf (pdfLaTeX)
% =====================================================================
\documentclass[12pt,a4paper]{report}

% ---------- Encodage & langue ----------
\usepackage[utf8]{inputenc}
\usepackage[T1]{fontenc}
\usepackage[french]{babel}
\usepackage{lmodern}

% ---------- Mise en page ----------
\usepackage[left=2.5cm,right=2.5cm,top=2.5cm,bottom=2.5cm]{geometry}
\usepackage{setspace}
\onehalfspacing
\usepackage{parskip}

% ---------- Mathématiques & symboles ----------
\usepackage{amsmath}
\usepackage{amsfonts}
\usepackage{amssymb}
\usepackage{textcomp}
\usepackage{pifont}

% ---------- Graphiques & tableaux ----------
\usepackage{graphicx}
\usepackage{float}
\usepackage{array}
\usepackage{booktabs}
\usepackage{longtable}
\usepackage{multirow}
\usepackage{xcolor}
\usepackage{colortbl}

% ---------- Listes ----------
\usepackage{enumitem}

% ---------- En-têtes / pieds de page ----------
\usepackage{fancyhdr}
\usepackage{titlesec}

% ---------- Liens ----------
\usepackage{hyperref}
\usepackage{csquotes}
\hypersetup{
    colorlinks=true,
    linkcolor=black,
    citecolor=black,
    filecolor=magenta,
    urlcolor=blue,
    pdftitle={Rapport de Stage - AI Network Copilot},
    pdfauthor={AI Network Copilot}
}
\urlstyle{same}

% ---------- Couleurs du document ----------
\definecolor{bleutt}{RGB}{0,84,159}
\definecolor{grisclair}{RGB}{240,240,240}
\definecolor{grismoyen}{RGB}{120,120,120}

% ---------- Style des titres de chapitre ----------
\titleformat{\chapter}[display]
  {\normalfont\huge\bfseries\color{bleutt}}
  {\chaptertitlename\ \thechapter}{20pt}{\Huge}
\titlespacing*{\chapter}{0pt}{10pt}{30pt}

% ---------- En-têtes ----------
\pagestyle{fancy}
\fancyhf{}
\rhead{\small\textit{AI Network Copilot}}
\lhead{\small\textit{Rapport de Stage}}
\rfoot{\thepage}
\renewcommand{\headrulewidth}{0.4pt}
\renewcommand{\footrulewidth}{0.4pt}

\fancypagestyle{plain}{%
  \fancyhf{}
  \rfoot{\thepage}
  \renewcommand{\headrulewidth}{0pt}
}

% ---------- Raccourcis ----------
\newcommand{\checkitem}{\ding{51}}      % coche
\newcommand{\diamonditem}{\ding{117}}   % losange
\newcommand{\redcircle}{\textcolor{red}{\ding{108}}}

% Boîte de remplacement pour une figure (compile sans fichier image)
\newcommand{\figplaceholder}[2][8cm]{%
  \begin{center}
  \setlength{\fboxsep}{6pt}%
  \fbox{\begin{minipage}[c][#1][c]{0.86\textwidth}
  \centering\itshape\color{grismoyen} #2
  \end{minipage}}
  \end{center}}

% =====================================================================
\begin{document}

% =====================================================================
%  PAGE DE GARDE
% =====================================================================
\begin{titlepage}
\thispagestyle{empty}
\begin{flushright}
\small Réf : \dotfill
\end{flushright}
\begin{flushright}
\small AU : 2025--2026
\end{flushright}

\vspace{0.4cm}
\begin{center}
{\large\textbf{Université de Sousse}}\\[0.2cm]
{\large\textbf{École Nationale d'Ingénieurs de Sousse}}\\[0.6cm]

\rule{\textwidth}{1pt}
\vspace{0.8cm}

{\Huge\bfseries\color{bleutt} Rapport de Stage d'été}\\[0.6cm]
{\large $[\ ]$ Initiation \qquad $[\times]$ Ingénieur}\\[1.2cm]

\rule{0.6\textwidth}{0.6pt}
\vspace{1cm}

{\large\bfseries Sujet :}\\[0.3cm]
{\LARGE\bfseries\color{bleutt} AI Network Copilot}\\[0.4cm]
{\large\itshape Système intelligent de supervision réseau et de détection\\ d'anomalies par intelligence artificielle}\\[1.4cm]

\rule{0.6\textwidth}{0.6pt}
\vspace{0.8cm}

\begin{minipage}{0.9\textwidth}
\centering
\large
\textbf{Réalisé par :}\\
[Nom et Prénom de l'étudiant(e)]\\[0.5cm]

\textbf{Filière :}\\
Génie Télécommunications et Réseaux\\[0.5cm]

\textbf{Stage effectué au sein de la Société (ou de l'Entreprise) :}\\
{[Nom de l'entreprise]}\\
Adresse : [Adresse de l'entreprise]\\[0.5cm]

\textbf{Encadrant :} [Nom de l'encadrant]\\[0.5cm]

\textbf{Période du stage :} [01/07/2026] -- [01/08/2026]
\end{minipage}

\end{center}
\end{titlepage}

% =====================================================================
%  REMERCIEMENTS
% =====================================================================
\chapter*{Remerciements}
\addcontentsline{toc}{chapter}{Remerciements}
\markboth{Remerciements}{Remerciements}

Je tiens à exprimer mes sincères remerciements à toute l'équipe de
[Nom de l'entreprise] pour son accueil chaleureux, son encadrement et son
soutien tout au long de mon stage d'ingénieur au sein de cette structure.
Leur professionnalisme et leur bienveillance ont grandement contribué à faire
de cette expérience une période enrichissante et formatrice.

Je suis particulièrement reconnaissant(e) envers mon encadrant,
[Nom de l'encadrant], pour sa disponibilité constante, ses précieux conseils et
sa patience. Grâce à ses orientations avisées, j'ai pu mener à bien mes missions
et approfondir mes connaissances dans le domaine de la supervision des réseaux,
de l'automatisation et de l'intelligence artificielle appliquée à la détection
d'anomalies.

Je tiens également à remercier l'ensemble des collaborateurs pour leur aide,
leurs conseils et leur esprit d'équipe, qui ont favorisé mon apprentissage et
mon épanouissement professionnel. Leur collaboration et leur expertise ont créé
un environnement de travail stimulant et convivial.

Enfin, j'exprime ma gratitude envers le corps enseignant de l'École Nationale
d'Ingénieurs de Sousse pour la qualité de la formation dispensée, ainsi qu'envers
toutes les personnes qui, de près ou de loin, ont contribué à la réussite de ce
projet. Cette expérience m'a permis d'appliquer les connaissances acquises au
cours de mes études, de découvrir le fonctionnement d'un projet technique complet
et d'améliorer mes compétences techniques et relationnelles.

% =====================================================================
%  SOMMAIRE / TABLE DES MATIÈRES
% =====================================================================
\renewcommand{\contentsname}{Sommaire}
\tableofcontents
\thispagestyle{fancy}

% =====================================================================
%  LISTE DES FIGURES
% =====================================================================
\renewcommand{\listfigurename}{Liste des figures}
\listoffigures
\addcontentsline{toc}{chapter}{Liste des figures}

% =====================================================================
%  FICHE-SYNTHÈSE DU STAGE
% =====================================================================
\chapter*{Fiche-synthèse du stage}
\addcontentsline{toc}{chapter}{Fiche-synthèse du stage}
\markboth{Fiche-synthèse du stage}{Fiche-synthèse du stage}

\renewcommand{\arraystretch}{1.4}
\begin{longtable}{|>{\raggedright\arraybackslash}p{6.2cm}|>{\raggedright\arraybackslash}p{8.5cm}|}
\hline
\rowcolor{grisclair}
\textbf{Rubrique} & \textbf{Renseignement} \\
\hline
Nom de l'entreprise & [Nom de l'entreprise] \\
\hline
Ville & [Ville] \\
\hline
Nom de l'encadrant industriel & [Nom de l'encadrant] \\
\hline
Date début et date fin du stage & 01/07/2026 -- 01/08/2026 \\
\hline
Nombre de jours effectifs passés dans les locaux de l'entreprise & 20 jours \\
\hline
Les principales tâches réalisées durant le stage (maximum 3) &
\begin{minipage}[t]{\linewidth}
\begin{itemize}[leftmargin=*,nosep]
\item Concevoir et automatiser une collecte temps réel des métriques réseau (latence, taux d'erreurs, trafic) et découvrir les équipements du réseau.
\item Développer un moteur de détection d'anomalies hybride (règles métier + Machine Learning) avec identification de la cause probable.
\item Réaliser une API (FastAPI) et un tableau de bord temps réel (Streamlit) pour la visualisation et l'alerte.
\end{itemize}
\end{minipage} \\
\hline
Le stage contient-il une période d'observation ? & Non \\
\hline
Le stage contient-il une partie conception ? Si oui, combien de jours ? & Oui -- 5 jours \\
\hline
Le stage contient-il une partie développement ? Si oui, combien de jours ? & Oui -- 12 jours \\
\hline
Prérequis en termes de compétences techniques &
Connaissances en réseaux informatiques, programmation Python, notions de Machine Learning \\
\hline
Les outils et logiciels utilisés durant le stage &
Python, FastAPI, Streamlit, SQLite, scikit-learn, pandas, Plotly, Docker, Git, VS Code \\
\hline
Le matériel (autre que les PCs) utilisé & Non \\
\hline
Méthodologies utilisées durant le stage & Approche incrémentale et itérative, travail collaboratif (contrat d'API), tests et validation \\
\hline
Deux principales compétences comportementales acquises & La rigueur / La résolution de problèmes \\
\hline
Deux principales compétences techniques acquises &
Développement d'API et de tableaux de bord temps réel, détection d'anomalies par intelligence artificielle \\
\hline
\end{longtable}
\renewcommand{\arraystretch}{1.0}

% =====================================================================
%  INTRODUCTION GÉNÉRALE
% =====================================================================
\chapter*{Introduction générale}
\addcontentsline{toc}{chapter}{Introduction générale}
\markboth{Introduction générale}{Introduction générale}

Au cours des dernières années, la complexité et la criticité des infrastructures
réseau ont considérablement augmenté, rendant leur gestion et leur supervision de
plus en plus exigeantes. Les interruptions de service, les dégradations de
performance et les incidents réseau représentent aujourd'hui des menaces majeures
pouvant engendrer des pertes économiques considérables et compromettre la
continuité des activités. Selon plusieurs études, une indisponibilité réseau de
seulement quelques minutes peut coûter plusieurs milliers d'euros à une entreprise
et affecter directement la productivité des utilisateurs.

Face à cette réalité, il est devenu impératif de mettre en place des solutions de
surveillance intelligente et proactive, capables non seulement de détecter les
incidents mais aussi d'en identifier la cause probable et de guider l'administrateur
vers la résolution, idéalement avant que l'incident n'affecte les utilisateurs.

Mon projet de stage s'inscrit dans cette problématique. En tant qu'étudiant(e) en
génie des télécommunications et réseaux, j'ai eu l'opportunité de concevoir et de
développer \textbf{AI Network Copilot}, un système intelligent de monitoring et de
maintenance prédictive. Ce système surveille en continu une infrastructure réseau,
découvre automatiquement les équipements connectés, collecte leurs métriques
(latence, taux d'erreurs, trafic), puis détecte les anomalies et en identifie la
cause au moyen d'un moteur hybride combinant des règles métier et un modèle de
Machine Learning.

La méthodologie employée repose sur l'intégration d'outils modernes de collecte et
de visualisation, couplés à des algorithmes d'apprentissage automatique
(\textit{Isolation Forest} pour la détection d'anomalies non supervisée) afin de
transformer des données brutes en décisions exploitables. L'ensemble de la
plateforme a été développé en Python, autour d'une API REST (FastAPI), d'une base
de données SQLite et d'un tableau de bord temps réel (Streamlit) enrichi de
graphiques interactifs (Plotly) et de communications temps réel (WebSocket).

Dans ce rapport, je commencerai par présenter l'organisme d'accueil et le contexte
de mon stage. J'exposerai ensuite la problématique de la supervision proactive et
l'intérêt d'utiliser l'intelligence artificielle pour détecter et expliquer les
anomalies. Puis je détaillerai la conception de la solution et la méthodologie
adoptée, de la collecte des métriques jusqu'à la génération des alertes. Enfin, je
présenterai la réalisation technique et les interfaces développées, avant de
conclure par un bilan de cette expérience enrichissante.

\vspace{0.5cm}
\noindent\textbf{Mots clés :} Supervision réseau / Découverte d'équipements /
Détection d'anomalies / Machine Learning / Isolation Forest / Maintenance prédictive /
Intelligence artificielle / FastAPI / Streamlit / Automatisation.

\setcounter{page}{1}
\pagenumbering{arabic}

% =====================================================================
%  CHAPITRE 1
% =====================================================================
\chapter{Organisme d'accueil et problématique traitée}

\section{Introduction}
Dans ce premier chapitre, nous explorerons l'environnement de travail dans lequel
s'est déroulé mon stage. Nous présenterons l'organisme d'accueil, son historique,
ses activités et sa structure organisationnelle. Nous exposerons ensuite la
problématique traitée et le travail demandé, la démarche méthodologique adoptée,
ainsi que le chronogramme des tâches menées durant le stage.

\section{Présentation de l'entreprise}

\subsection{Présentation générale}
[Nom de l'entreprise] est une structure évoluant dans le domaine des technologies
de l'information et des télécommunications. Elle propose à ses clients une gamme de
services couvrant l'ingénierie réseau, la supervision d'infrastructures, la
sécurité des systèmes et le développement de solutions logicielles. Son objectif
est d'accompagner les organisations dans la modernisation et la fiabilisation de
leurs infrastructures numériques.

\subsection{Historique}
[Nom de l'entreprise] a été fondée en [année de création]. Depuis sa création,
elle n'a cessé d'élargir son domaine d'expertise et d'améliorer la qualité de ses
services afin de répondre à l'évolution rapide des besoins technologiques de ses
clients. Elle s'est progressivement spécialisée dans les solutions de supervision
et d'automatisation, en intégrant les technologies émergentes telles que
l'intelligence artificielle et l'analyse de données.

\subsection{Activités}
L'entreprise joue un rôle important à travers ses services qui touchent de
nombreux aspects des systèmes d'information de ses clients. Elle est notamment
responsable de la conception, du déploiement et de la maintenance
d'infrastructures réseau, ainsi que de la mise en place de solutions de
monitoring et de sécurité. Elle offre des services aussi bien aux entreprises
qu'aux particuliers pour tout ce qui concerne la connectivité, la disponibilité et
la protection des systèmes.

\subsection{L'identité visuelle}
Le logo de [Nom de l'entreprise] reflète les valeurs de l'organisation :
innovation, fiabilité et proximité. La figure suivante présente l'identité
visuelle de l'entreprise.

\begin{figure}[H]
\figplaceholder[4cm]{Figure~1 : Logo de [Nom de l'entreprise]\\(insérer le logo de l'entreprise)}
\caption{Logo de l'entreprise}
\label{fig:logo}
\end{figure}

\subsection{Organisation générale de l'entreprise}
L'organisation de [Nom de l'entreprise] est modélisée par l'organigramme présenté
dans la figure ci-dessous. Elle se compose de plusieurs départements
(direction générale, département technique, département commercial,
support et administration) travaillant de manière coordonnée.

\begin{figure}[H]
\figplaceholder[5cm]{Figure~2 : Organigramme de l'entreprise\\(insérer l'organigramme)}
\caption{Organisation générale de l'entreprise}
\label{fig:organigramme}
\end{figure}

\section{Problématique traitée et travail demandé}
Dans les environnements informatiques actuels, la disponibilité et la fiabilité
des réseaux sont essentielles. Les infrastructures, composées d'un nombre croissant
d'équipements hétérogènes (postes de travail, serveurs, équipements réseau,
terminaux mobiles), doivent être surveillées en temps réel pour prévenir toute
interruption. La détection réactive des anomalies, c'est-à-dire après leur
apparition, entraîne souvent des temps d'arrêt, des pertes de productivité et des
risques accrus pour la sécurité. Une approche proactive, capable non seulement de
détecter mais aussi d'\emph{expliquer} les anomalies, est donc indispensable pour
garantir performance et protection.

Un autre défi majeur réside dans le fait que les administrateurs sont fréquemment
submergés par un grand nombre d'alertes brutes, sans indication claire sur la
cause du problème ni sur la marche à suivre pour le résoudre. Il en résulte une
perte de temps précieuse dans le diagnostic manuel.

Le travail demandé consistait donc à concevoir et à développer un système de
supervision proactive et intelligente d'une infrastructure réseau, baptisé
\textbf{AI Network Copilot}, permettant de :
\begin{itemize}[leftmargin=*]
\item découvrir automatiquement les équipements présents sur le réseau et les
identifier (adresse IP, adresse MAC, fabricant, type d'appareil) ;
\item collecter en continu les métriques réseau clés (latence, taux d'erreurs,
trafic) ;
\item détecter les anomalies à l'aide d'un moteur hybride combinant des règles
métier et un modèle d'intelligence artificielle ;
\item identifier automatiquement la \emph{cause probable} de chaque anomalie et
proposer un conseil de diagnostic à l'administrateur ;
\item visualiser en temps réel l'état du réseau et les alertes via un tableau de
bord interactif, et exposer l'ensemble des données à travers une API REST.
\end{itemize}

\section{Démarche -- Méthodologie utilisée}
Pour répondre à cette problématique, la méthodologie suivante a été adoptée :
\begin{itemize}[leftmargin=*]
\item \textbf{\diamonditem\ Analyse des besoins et des technologies existantes :}
étude des solutions de monitoring du marché, identification de leurs limites
(seuils statiques, absence d'explication des causes) et choix d'outils adaptés
(Python, FastAPI, Streamlit, scikit-learn) pour collecter, analyser et visualiser
les métriques.
\item \textbf{\diamonditem\ Conception de l'architecture du système :}
élaboration d'une architecture modulaire intégrant la collecte des métriques et la
découverte des équipements, la détection d'anomalies par IA, l'identification des
causes, la visualisation temps réel et la génération d'alertes.
\item \textbf{\diamonditem\ Développement et implémentation :} mise en place du
collecteur (\textit{ping}, découverte ARP), de la base de données SQLite, du
moteur de détection (règles + \textit{Isolation Forest}), de l'API FastAPI et du
tableau de bord Streamlit.
\item \textbf{\diamonditem\ Intégration temps réel :} mise en place de canaux
WebSocket pour la diffusion en direct des métriques, des alertes et de la
topologie vers le tableau de bord.
\item \textbf{\diamonditem\ Tests et validation :} évaluation de la précision de la
détection, contrôle du bon fonctionnement des différents services, empaquetage de
l'application et ajustement des seuils pour garantir la fiabilité de la solution.
\end{itemize}

\section{Chronogramme des tâches menées durant le stage}
La figure suivante présente le chronogramme (diagramme de Gantt) des principales
tâches réalisées durant la période du stage.

\begin{figure}[H]
\figplaceholder[5.5cm]{Figure~3 : Chronogramme (diagramme de Gantt) des tâches\\(insérer le planning du stage)}
\caption{Chronogramme des tâches menées durant le stage}
\label{fig:gantt}
\end{figure}

\noindent\textbf{Explication :}

\noindent\textbf{Tâche 1 -- Recherche et auto-formation :}
\begin{itemize}[leftmargin=*,nosep]
\item Recherche approfondie sur les technologies existantes de détection
d'anomalies réseau.
\item Étude des bibliothèques Python (scikit-learn, pandas, NumPy) pour le
développement du modèle d'IA.
\item Prise en main des frameworks FastAPI (API) et Streamlit (tableau de bord).
\item Documentation des meilleures pratiques et des approches de projets
similaires.
\end{itemize}

\noindent\textbf{Tâche 2 -- Conception du système :}
\begin{itemize}[leftmargin=*,nosep]
\item Définition de l'architecture globale et du schéma de la base de données.
\item Établissement d'un contrat d'API entre le module de collecte/détection et le
tableau de bord.
\item Conception du diagramme de flux du système de supervision proactive.
\end{itemize}

\noindent\textbf{Tâche 3 -- Collecte et découverte :}
\begin{itemize}[leftmargin=*,nosep]
\item Développement du collecteur de métriques (\textit{ping}) et de la découverte
automatique des équipements (balayage de sous-réseau, cache ARP, sondes TCP).
\item Identification des équipements (adresse MAC, fabricant via OUI, nom d'hôte,
type d'appareil) et stockage dans la base SQLite.
\end{itemize}

\noindent\textbf{Tâche 4 -- Détection, API et visualisation :}
\begin{itemize}[leftmargin=*,nosep]
\item Implémentation du moteur de détection hybride (règles + \textit{Isolation
Forest}) avec identification de la cause et conseil associé.
\item Développement de l'API REST (FastAPI) et du tableau de bord temps réel
(Streamlit) avec graphiques Plotly et diffusion WebSocket.
\item Empaquetage de l'application et validation de l'ensemble de la chaîne.
\end{itemize}

\section{Conclusion : principales observations et constatations}
Ce chapitre a permis de présenter l'entreprise d'accueil, de définir la
problématique traitée et de détailler la démarche méthodologique ainsi que les
principales tâches effectuées durant le stage. La principale constatation est que
la supervision réactive traditionnelle est insuffisante face à la complexité
croissante des réseaux : il est nécessaire de disposer d'un système capable non
seulement de détecter les anomalies, mais aussi d'en expliquer la cause pour
accélérer le diagnostic. C'est précisément l'objectif d'\textbf{AI Network
Copilot}, dont la conception est détaillée dans le chapitre suivant.

% =====================================================================
%  CHAPITRE 2
% =====================================================================
\chapter{Analyse des besoins et conception}

\section{Introduction}
Dans ce chapitre, nous aborderons l'étude de l'existant, la solution proposée,
ainsi que la méthodologie et le processus retenus pour le système de supervision
proactive \textbf{AI Network Copilot}. Cette étape est cruciale pour définir les
fonctionnalités nécessaires et pour planifier la manière dont le système sera
construit.

\section{Étude de l'existant}
Au cours des dernières années, les infrastructures réseau sont devenues de plus en
plus complexes, intégrant des postes de travail, des serveurs, des équipements
d'interconnexion et de nombreux terminaux mobiles. Dans beaucoup d'organisations,
la surveillance repose encore sur des solutions \textbf{réactives}, où les
incidents ne sont traités qu'après leur apparition. Cette approche entraîne des
interruptions de service coûteuses, une perte de productivité et parfois des
failles de sécurité exploitables.

Bien que certaines plateformes de monitoring classiques (comme \textit{Nagios},
\textit{Zabbix} ou \textit{PRTG}) permettent de visualiser l'état des ressources,
elles présentent plusieurs limites :
\begin{itemize}[leftmargin=*]
\item \textbf{Alertes basées sur des seuils statiques :} ces seuils, fixés
manuellement, génèrent soit trop de fausses alertes, soit passent à côté de
dégradations progressives des performances.
\item \textbf{Absence d'explication :} elles signalent qu'un problème existe, mais
n'indiquent ni sa cause probable ni la marche à suivre pour le résoudre.
\item \textbf{Configuration lourde :} leur déploiement et leur paramétrage
requièrent souvent une expertise importante et un effort de maintenance continu.
\item \textbf{Faible capacité prédictive :} la plupart ne s'appuient pas sur des
techniques d'apprentissage automatique pour anticiper les incidents.
\end{itemize}

Face à ces limites, l'adoption d'une solution \textbf{intelligente, légère et
explicative}, basée sur l'intelligence artificielle, devient pertinente.
L'objectif est de détecter les anomalies, d'en identifier automatiquement la
cause, d'optimiser la disponibilité des services et de guider l'administrateur
dans la résolution, tout en réduisant les coûts liés aux arrêts imprévus.

\section{Solution}
Pour surmonter les limites des approches classiques de supervision réseau, nous
avons conçu une solution innovante de surveillance proactive baptisée
\textbf{AI Network Copilot}, s'appuyant sur l'intelligence artificielle et sur des
outils logiciels modernes et légers. L'objectif est de détecter et d'expliquer les
incidents avant qu'ils n'affectent durablement les utilisateurs ou la
disponibilité des services.

Notre solution repose sur \textbf{quatre piliers principaux} :

\subsection*{\diamonditem\ Collecte des données et découverte du réseau}
La collecte des métriques réseau est réalisée par un \textbf{collecteur} qui
mesure, pour chaque équipement, la latence, le taux d'erreurs et le trafic. Le
système intègre également un mécanisme de \textbf{découverte automatique} des
équipements du réseau : balayage de sous-réseau par \textit{ping}, lecture du
cache ARP/voisinage (\textit{Get-NetNeighbor}), balayage TCP et \textit{ping} de
diffusion. Chaque équipement est ensuite identifié (adresse IP, adresse MAC,
fabricant déduit de l'OUI, nom d'hôte, type d'appareil). Ces données sont
centralisées et historisées dans une base de données \textbf{SQLite}.

\subsection*{\diamonditem\ Détection d'anomalies par intelligence artificielle}
Les données collectées alimentent un \textbf{moteur de détection hybride} à deux
couches :
\begin{itemize}[leftmargin=*,nosep]
\item une couche de \textbf{règles métier} qui détecte les dépassements de seuils
(latence, taux d'erreurs) et identifie la cause probable ;
\item une couche de \textbf{Machine Learning} fondée sur l'algorithme
\textit{Isolation Forest} de la bibliothèque scikit-learn. Ce modèle
non supervisé apprend le comportement normal du réseau et détecte les mesures
statistiquement anormales, y compris celles que des seuils fixes ne captureraient
pas.
\end{itemize}
Une mesure est considérée comme anormale si \emph{l'une} des deux couches se
déclenche. Le système identifie alors la \textbf{cause probable} et fournit un
\textbf{conseil de diagnostic} ainsi qu'un \textbf{score de confiance}.

\subsection*{\diamonditem\ Visualisation en temps réel}
Pour offrir une vue globale et exploitable de l'infrastructure, un tableau de bord
interactif a été développé avec \textbf{Streamlit} et \textbf{Plotly}. Il permet de
visualiser en temps réel les métriques de performance, les anomalies détectées, la
liste des équipements et une \textbf{vue topologique} du réseau. La mise à jour en
direct est assurée par des canaux \textbf{WebSocket}, avec repli automatique sur des
requêtes HTTP.

\subsection*{\diamonditem\ Alertes et identification des causes}
Lorsqu'une anomalie est détectée, le système génère une \textbf{alerte enrichie}
comprenant l'équipement concerné, la cause probable, un conseil de diagnostic et un
score de confiance. Cette approche transforme un simple signalement en une
information directement actionnable par l'administrateur.

Cette solution, combinant un collecteur léger, une base SQLite, un moteur de
détection hybride, une API REST (FastAPI) et un tableau de bord temps réel
(Streamlit), permet de faire passer la supervision réseau d'une approche
\textbf{réactive} à une approche \textbf{prédictive, proactive et explicative},
garantissant une meilleure disponibilité et une réduction du temps de diagnostic.

\section{Méthodologie et processus}
Cette partie aborde l'approche utilisée pour atteindre les objectifs du projet et
examine de plus près la manière dont le système est mis en œuvre.

\subsection{Architecture du système}
Le système est organisé en modules découplés qui communiquent via la base de
données et l'API :
\begin{itemize}[leftmargin=*,nosep]
\item \textbf{Collector} : découverte des équipements et collecte des métriques.
\item \textbf{Base de données (SQLite)} : stockage des métriques et des propriétés
des équipements.
\item \textbf{Detection} : règles métier + modèle \textit{Isolation Forest} +
identification des causes.
\item \textbf{API (FastAPI)} : exposition des métriques, alertes, hôtes,
équipements et topologie (REST + WebSocket).
\item \textbf{Dashboard (Streamlit)} : visualisation temps réel et interaction
utilisateur.
\end{itemize}

\begin{figure}[H]
\figplaceholder[6cm]{Figure~4 : Architecture modulaire d'AI Network Copilot\\Collector $\rightarrow$ Base SQLite $\rightarrow$ Detection (Règles + ML) $\rightarrow$ API FastAPI $\rightarrow$ Dashboard Streamlit\\(insérer le schéma d'architecture)}
\caption{Architecture globale du système AI Network Copilot}
\label{fig:archi}
\end{figure}

\subsection{Diagramme de flux du système}
La figure suivante présente le processus complet du système de surveillance
proactive grâce à l'IA.

\begin{figure}[H]
\figplaceholder[6cm]{Figure~5 : Schéma de processus complet du système de surveillance proactive grâce à l'IA\\(Data Collection $\rightarrow$ Feature Extraction $\rightarrow$ Détection Règles + Isolation Forest $\rightarrow$ Identification de la cause $\rightarrow$ Visualisation $\rightarrow$ Alerting)}
\caption{Schéma de processus complet du système de surveillance proactive grâce à l'IA}
\label{fig:flux}
\end{figure}

\noindent\textbf{Data Collection :}
Le système assure une collecte continue et centralisée des métriques réseau
provenant de chaque équipement découvert. Ces données incluent des indicateurs
critiques tels que la latence, le taux d'erreurs et le trafic. Grâce au collecteur
et à la base SQLite, la collecte est automatisée, fiable et continue, garantissant
une base de données exploitable pour l'analyse en temps réel et la détection
proactive d'anomalies.

\noindent\textbf{Feature Extraction :}
Sélection des indicateurs réseau les plus pertinents pour la détection
d'anomalies (latence, taux d'erreurs, trafic), afin d'alimenter le modèle d'IA
avec des données significatives et de maximiser la précision des prédictions.

\noindent\textbf{AI Model Inference (Isolation Forest + Règles) :}
Le moteur hybride analyse les indicateurs sélectionnés pour déterminer, en temps
réel, si le comportement observé est normal ou anormal. La couche règles capture
les dépassements de seuils explicites tandis que l'\textit{Isolation Forest}
capture les anomalies statistiques subtiles.

\noindent\textbf{Identification de la cause :}
En cas d'anomalie, le système déduit la cause probable (par ex. machine
inaccessible, perte de paquets, congestion, saturation de bande passante, panne
WAN, etc.) à partir de la combinaison des valeurs mesurées, et associe à chaque
cause un conseil de diagnostic.

\noindent\textbf{Visualisation :}
Le système utilise un tableau de bord Streamlit/Plotly pour afficher en temps réel
les résultats, avec des graphiques interactifs des indicateurs de santé du réseau,
facilitant le suivi et la prise de décision.

\noindent\textbf{Alerting \& actions :}
Lorsqu'une anomalie est détectée, le système produit une alerte enrichie
(cause + conseil + score de confiance) et la diffuse en temps réel vers le tableau
de bord via WebSocket, permettant une réaction rapide de l'administrateur.

\section{Conclusion}
Ce chapitre a permis de concevoir une solution adaptée en analysant les limites
des solutions existantes. La méthodologie suivie a structuré le système en modules
cohérents (collecte, détection, API, visualisation) et a préparé sa mise en œuvre,
en mettant l'accent sur une approche méthodique, progressive et explicative. Le
chapitre suivant détaille la réalisation technique de cette conception.

% =====================================================================
%  CHAPITRE 3
% =====================================================================
\chapter{Réalisation}

\section{Introduction}
Dans ce chapitre, je présenterai les choix technologiques et les environnements de
développement qui ont guidé la réalisation du système \textbf{AI Network Copilot}.
Je détaillerai ensuite les fonctionnalités implémentées, les interfaces
développées, ainsi que les tests effectués pour valider le bon fonctionnement du
système.

\section{Choix technologiques et environnements de développement}

\subsection{Les technologies et les outils utilisés}
Dans ce projet, un ensemble de technologies, de bibliothèques et d'outils logiciels
a été mis en œuvre pour concevoir un système léger et complet de supervision
réseau, assurant la collecte, l'analyse et la visualisation intelligente des
métriques, ainsi que la détection et l'explication proactive des anomalies. La
section suivante présente une description succincte de chaque outil et de son rôle.

\vspace{0.3cm}
\noindent\textbf{1. Python.}
Python est un langage de programmation de haut niveau, connu pour sa simplicité et
sa lisibilité. Grâce à son large écosystème de bibliothèques, il a été le langage
principal utilisé pour implémenter l'ensemble du système (collecte, détection, API
et tableau de bord).

\begin{figure}[H]
\figplaceholder[3cm]{Figure~6 : Logo Python}
\caption{Logo du langage Python}
\label{fig:python}
\end{figure}

\noindent\textbf{2. FastAPI.}
FastAPI est un framework web moderne et performant pour construire des API REST en
Python. Il a été utilisé pour exposer les endpoints du système
(\texttt{/metrics/latest}, \texttt{/alerts}, \texttt{/hosts}, \texttt{/devices},
\texttt{/topology}, \texttt{/ping/\{host\}}, \texttt{/health}) et pour la
communication temps réel via WebSocket. Il fournit également une documentation
interactive automatique (Swagger).

\begin{figure}[H]
\figplaceholder[3cm]{Figure~7 : Logo FastAPI}
\caption{Logo du framework FastAPI}
\label{fig:fastapi}
\end{figure}

\noindent\textbf{3. Streamlit.}
Streamlit est un framework Python permettant de créer rapidement des applications
web de visualisation de données. Il a été utilisé pour développer le tableau de
bord interactif temps réel du système.

\begin{figure}[H]
\figplaceholder[3cm]{Figure~8 : Logo Streamlit}
\caption{Logo du framework Streamlit}
\label{fig:streamlit}
\end{figure}

\noindent\textbf{4. SQLite.}
SQLite est un système de gestion de base de données relationnelle léger, sans
serveur. Il stocke les métriques réseau (table \texttt{network\_metrics}) et les
propriétés des équipements (table \texttt{device\_properties}), assurant
l'historisation et l'interrogation efficace des données.

\begin{figure}[H]
\figplaceholder[3cm]{Figure~9 : Logo SQLite}
\caption{Logo de SQLite}
\label{fig:sqlite}
\end{figure}

\noindent\textbf{5. scikit-learn (Isolation Forest).}
scikit-learn est une bibliothèque Python de référence pour le Machine Learning.
Dans ce projet, elle permet la conception et l'entraînement du modèle de détection
d'anomalies non supervisé fondé sur l'algorithme \textit{Isolation Forest}, capable
d'identifier les comportements statistiquement anormaux du réseau.

\begin{figure}[H]
\figplaceholder[3cm]{Figure~10 : Logo scikit-learn}
\caption{Logo de scikit-learn}
\label{fig:sklearn}
\end{figure}

\noindent\textbf{6. Plotly.}
Plotly est une bibliothèque de visualisation interactive. Elle est utilisée dans le
tableau de bord pour tracer des courbes temporelles de latence, de trafic et de
taux d'erreurs, ainsi que pour la mise en évidence visuelle des anomalies.

\noindent\textbf{7. WebSocket.}
Le protocole WebSocket permet une communication bidirectionnelle en temps réel
entre le serveur (API) et le client (tableau de bord). Il est utilisé pour diffuser
en continu les mises à jour des métriques, des alertes et de la topologie.

\noindent\textbf{8. Bibliothèques Python complémentaires.}
\begin{itemize}[leftmargin=*,nosep]
\item \textbf{pandas} : manipulation et analyse des données ; structuration et
préparation des métriques avant leur passage au modèle d'IA.
\item \textbf{NumPy} : calcul numérique efficace sur des tableaux, utilisé lors du
traitement des données.
\item \textbf{joblib} : sérialisation et sauvegarde du modèle entraîné
(\texttt{model.pkl}) afin de le recharger rapidement en production.
\item \textbf{scapy} et \textbf{ping3} : découverte réseau (ARP) et mesure de la
latence par \textit{ping}.
\end{itemize}

\noindent\textbf{9. Docker.}
Docker permet de conteneuriser l'API et le tableau de bord afin de faciliter leur
déploiement de manière reproductible, indépendamment de l'environnement hôte.

\noindent\textbf{10. Git et Visual Studio Code.}
Git a été utilisé pour le versionnement du code et le travail collaboratif. Visual
Studio Code, éditeur développé par Microsoft, a servi à l'écriture, à l'édition et
au débogage des scripts Python du projet.

\noindent\textbf{11. PyInstaller.}
PyInstaller a été utilisé pour empaqueter l'application sous forme d'un exécutable
Windows autonome (\texttt{AI-Network-Copilot.exe}), lancé par un \textit{launcher}
qui orchestre les trois services (collecteur, API, tableau de bord).

\subsection{Le matériel utilisé}
Le développement et les tests de ce projet ont été réalisés sur un ordinateur
portable présentant les spécifications suivantes.

\begin{table}[H]
\centering
\renewcommand{\arraystretch}{1.3}
\begin{tabular}{|>{\raggedright\arraybackslash}p{4cm}|>{\raggedright\arraybackslash}p{6cm}|}
\hline
\rowcolor{grisclair}
\multicolumn{2}{|c|}{\textbf{PC1 -- Caractéristiques}} \\
\hline
Processeur & Intel Core i5 \\
\hline
RAM & 16 Go \\
\hline
Système d'exploitation & Windows 11 \\
\hline
Stockage & SSD \\
\hline
\end{tabular}
\caption{Caractéristiques du PC de développement}
\label{tab:pc}
\end{table}

\section{Les fonctionnalités réalisées}
Le système développé intègre les fonctionnalités suivantes :

\begin{itemize}[leftmargin=*]
\item \textbf{\checkitem\ Découverte automatique du réseau :} le collecteur
détecte automatiquement les équipements présents sur le(s) sous-réseau(x) via un
balayage \textit{ping}, la lecture du cache ARP/voisinage
(\textit{Get-NetNeighbor}), des sondes TCP et un \textit{ping} de diffusion. Chaque
équipement est identifié (adresse IP, adresse MAC, fabricant via l'OUI, nom d'hôte,
type d'appareil).

\item \textbf{\checkitem\ Collecte et prétraitement des données :} les métriques
(latence, taux d'erreurs, trafic) sont mesurées puis nettoyées et structurées avant
d'être stockées dans la base SQLite, garantissant la qualité et la cohérence des
données utilisées par le système.

\item \textbf{\checkitem\ Extraction des caractéristiques :} les indicateurs les
plus pertinents pour la détection d'anomalies (latence, taux d'erreurs, trafic)
sont sélectionnés afin d'alimenter le modèle d'IA avec des données significatives.

\item \textbf{\checkitem\ Développement et entraînement du modèle
(Isolation Forest) :} scikit-learn a été utilisé pour concevoir et entraîner un
modèle non supervisé capable de détecter les comportements anormaux. Le modèle est
sauvegardé (\texttt{model.pkl}) et rechargé en cache pour des inférences rapides,
individuelles ou par lot.

\item \textbf{\checkitem\ Détection hybride et identification des causes :} un
moteur combine des règles métier (seuils de latence et de taux d'erreurs) et le
modèle ML. Une mesure est anormale si l'une des deux couches se déclenche. Le
système déduit alors la cause probable (machine inaccessible, perte de paquets,
congestion, saturation de bande passante, panne WAN, problème DNS, etc.) et fournit
un conseil de diagnostic ainsi qu'un score de confiance.

\item \textbf{\checkitem\ API REST et temps réel :} l'API FastAPI expose les
métriques, alertes, hôtes, équipements et la topologie du réseau. Des canaux
WebSocket (\texttt{/ws/metrics}, \texttt{/ws/alerts}, \texttt{/ws/topology})
diffusent les mises à jour en temps réel.

\item \textbf{\checkitem\ Visualisation des résultats (Dashboard) :} le tableau de
bord Streamlit/Plotly affiche en temps réel les métriques, les alertes enrichies,
la liste des équipements et une vue topologique du réseau, facilitant le suivi
continu et la prise de décision rapide. Une fonction de \textit{ping} en direct
permet de tester un hôte à la demande.

\item \textbf{\checkitem\ Empaquetage et déploiement :} l'application est
distribuée sous forme d'exécutable Windows autonome et peut également être déployée
via Docker. Un \textit{launcher} orchestre le démarrage coordonné des trois
services.

\item \textbf{\checkitem\ Supervision proactive et explicative :} la combinaison
de toutes les fonctionnalités précédentes forme un système complet capable de
détecter, d'expliquer et de signaler les anomalies en temps réel, assurant une
meilleure disponibilité et un diagnostic accéléré.
\end{itemize}

\section{Les interfaces et le prototype}

\subsection{Interface principale du tableau de bord}
La figure ci-dessous présente la vue principale du tableau de bord temps réel, avec
les indicateurs clés (KPIs), l'état global du réseau et les courbes de performance.

\begin{figure}[H]
\figplaceholder[6cm]{Figure~11 : Vue principale du tableau de bord AI Network Copilot\\(insérer la capture d'écran du dashboard Streamlit)}
\caption{Vue principale du tableau de bord temps réel}
\label{fig:dashboard}
\end{figure}

\subsection{Visualisation temps réel des métriques}
Les graphiques interactifs (Plotly) affichent l'évolution de la latence, du trafic
et du taux d'erreurs, avec mise en évidence des points d'anomalie.

\begin{figure}[H]
\figplaceholder[6cm]{Figure~12 : Courbes temporelles de latence / trafic / taux d'erreurs\\(insérer la capture d'écran des graphiques)}
\caption{Visualisation temps réel des métriques réseau}
\label{fig:metrics}
\end{figure}

\subsection{Liste et identification des équipements découverts}
Le système affiche la liste des équipements découverts avec leurs informations
(adresse IP, adresse MAC, fabricant, type d'appareil, dernière observation).

\begin{figure}[H]
\figplaceholder[6cm]{Figure~13 : Liste des équipements découverts sur le réseau\\(insérer la capture d'écran de la table des devices)}
\caption{Équipements découverts et identifiés}
\label{fig:devices}
\end{figure}

\subsection{Vue topologique du réseau}
La vue topologique représente les équipements sous forme de nœuds reliés par des
liens, colorés selon leur type et leur état.

\begin{figure}[H]
\figplaceholder[6cm]{Figure~14 : Vue topologique du réseau\\(insérer la capture d'écran de la topologie)}
\caption{Vue topologique du réseau}
\label{fig:topology}
\end{figure}

\subsection{Détection d'anomalies et identification des causes}
Les figures ci-dessous illustrent la détection d'une anomalie ainsi que l'alerte
enrichie générée par le système, comprenant la cause probable et le conseil de
diagnostic.

\begin{figure}[H]
\figplaceholder[6cm]{Figure~15 : Anomalie détectée avec identification de la cause et conseil\\(insérer la capture d'écran de l'alerte enrichie)}
\caption{Anomalie détectée avec cause probable et conseil de diagnostic}
\label{fig:alert}
\end{figure}

\subsection{Documentation de l'API}
L'API FastAPI fournit une documentation interactive (Swagger) accessible via
\texttt{/docs}, permettant de tester directement les endpoints.

\begin{figure}[H]
\figplaceholder[6cm]{Figure~16 : Documentation interactive de l'API (Swagger)\\(insérer la capture d'écran de /docs)}
\caption{Documentation interactive de l'API (Swagger)}
\label{fig:swagger}
\end{figure}

\section{Tests et validation}
Plusieurs tests ont été menés pour valider le bon fonctionnement de la solution :
\begin{itemize}[leftmargin=*,nosep]
\item \textbf{Tests unitaires du moteur de détection :} vérification, sur des cas
types (machine inaccessible, congestion, perte de paquets, panne WAN, etc.), de
l'exactitude de l'identification de la cause et de la cohérence du conseil associé.
\item \textbf{Tests d'intégration :} validation de la chaîne complète
collecteur $\rightarrow$ base SQLite $\rightarrow$ détection $\rightarrow$ API
$\rightarrow$ tableau de bord.
\item \textbf{Tests temps réel :} contrôle de la diffusion WebSocket et du repli
HTTP en cas d'indisponibilité.
\item \textbf{Tests de robustesse :} gestion des cas d'erreur (hôte injoignable,
réseau de type \textit{hotspot}, base vide) et repli sur des données simulées pour
garantir la continuité de l'affichage.
\end{itemize}
Ces tests ont confirmé que la solution détecte et signale rapidement les anomalies,
tout en fournissant une explication exploitable de leur cause probable.

% =====================================================================
%  CONCLUSION GÉNÉRALE
% =====================================================================
\chapter*{Conclusion générale}
\addcontentsline{toc}{chapter}{Conclusion générale}
\markboth{Conclusion générale}{Conclusion générale}

Au cours de ce projet, j'ai travaillé sur le développement d'\textbf{AI Network
Copilot}, un système de supervision réseau capable de découvrir automatiquement les
équipements connectés, de collecter leurs métriques, de détecter les comportements
anormaux et d'en identifier la cause probable. La problématique principale était
d'assurer une supervision \emph{proactive, fiable et explicative} du réseau afin de
réduire les interruptions, d'accélérer le diagnostic et de renforcer la sécurité.

La démarche méthodologique adoptée a impliqué plusieurs étapes clés : l'analyse des
besoins et des solutions existantes, la conception d'une architecture modulaire, la
collecte et le prétraitement des métriques, le développement d'un moteur de
détection hybride (règles métier + \textit{Isolation Forest}), ainsi que la
réalisation d'une API REST et d'un tableau de bord temps réel. L'approche a intégré
des outils et bibliothèques tels que Python, FastAPI, Streamlit, SQLite,
scikit-learn, pandas, Plotly et Docker.

Les fonctionnalités développées incluent la découverte et l'identification
automatiques des équipements, la détection en temps réel des anomalies, l'analyse
proactive des indicateurs clés (latence, taux d'erreurs, trafic), l'identification
de la cause probable avec conseil de diagnostic, l'affichage dynamique des
métriques, des alertes et de la topologie via un tableau de bord interactif, ainsi
que l'empaquetage de l'application pour un déploiement simple. Ces fonctionnalités
permettent de surveiller en continu l'état du réseau et d'anticiper les incidents.

La solution développée s'est révélée performante pour détecter, expliquer et
signaler rapidement les anomalies. Certaines améliorations pourraient être
envisagées à l'avenir, notamment l'extension du modèle à de nouveaux indicateurs,
l'ajout de seuils adaptatifs par équipement, l'intégration de notifications
externes (courriel, messagerie) et l'usage de techniques d'apprentissage plus
avancées pour renforcer encore la précision et la fiabilité du système.

Ce projet a été très enrichissant tant sur le plan technique que méthodologique.
J'ai acquis des compétences solides en conception et déploiement de systèmes de
supervision réseau, en développement d'API et de tableaux de bord temps réel, ainsi
qu'en Machine Learning appliqué à la détection d'anomalies. Sur le plan
organisationnel, il m'a permis de structurer un projet complexe de manière cohérente
et d'appréhender les défis liés à la mise en œuvre d'un système complet.

Cette expérience a renforcé ma motivation à poursuivre une carrière dans les
domaines de la supervision réseau, de la cybersécurité et du développement de
solutions intelligentes, et m'a offert de nouvelles perspectives pour mon avenir
professionnel.

% =====================================================================
%  RÉFÉRENCES
% =====================================================================
\chapter*{Références}
\addcontentsline{toc}{chapter}{Références}
\markboth{Références}{Références}

\begin{enumerate}[label={[\arabic*]}]
\item Python Software Foundation, \emph{Python Language Documentation}.
\url{https://www.python.org/} [consulté en 2026].
\item Sebastián Ramírez, \emph{FastAPI Documentation}.
\url{https://fastapi.tiangolo.com/} [consulté en 2026].
\item Streamlit Inc., \emph{Streamlit Documentation}.
\url{https://docs.streamlit.io/} [consulté en 2026].
\item SQLite Consortium, \emph{SQLite Documentation}.
\url{https://www.sqlite.org/docs.html} [consulté en 2026].
\item F. Pedregosa \emph{et al.}, \enquote{Scikit-learn: Machine Learning in
Python}, \emph{Journal of Machine Learning Research}, vol.~12, 2011.
\url{https://scikit-learn.org/} [consulté en 2026].
\item F. T. Liu, K. M. Ting, Z.-H. Zhou, \enquote{Isolation Forest},
\emph{IEEE International Conference on Data Mining (ICDM)}, 2008.
\item Plotly Technologies Inc., \emph{Plotly Python Graphing Library}.
\url{https://plotly.com/python/} [consulté en 2026].
\item Docker Inc., \emph{Docker Documentation}.
\url{https://docs.docker.com/} [consulté en 2026].
\item Microsoft, \emph{Visual Studio Code Documentation}.
\url{https://code.visualstudio.com/} [consulté en 2026].
\end{enumerate}

\end{document}



